\documentclass[11pt]{article}
\usepackage{amsmath,amssymb}
\usepackage[hidelinks]{hyperref}
\title{When a new subscriber changes the cost of already queued work}
\author{}\date{}
\begin{document}\maketitle

\paragraph{Scope.}
Two completion reports are already permitted and queued. They use
different senders and required receivers, but the senders can each have
an optional 60-byte probe addressed to a common endpoint. Reports are
173 bytes and carry an exact named output of pinned MRWS completion.
The tests use pinned SubT routing, native sender debits and callbacks
with deterministic path and pose seams. Report admission, capacity
envelopes and the right to defer a new subscription are fixture inputs.
No hostile event or field intervention is considered.

\paragraph{Rule.}
The \href{https://github.com/osrf/subt/blob/083c6a8acb0c25aa1a35e34517e628cd3b52ae77/subt-communication/subt_communication_broker/src/subt_communication_broker.cpp}{pinned broker}
looks up endpoint bindings at dispatch, not at initial enqueue. Bind
adds a receiver to the end of the native endpoint vector. For required
attempt \(j\) of sender \(s\), let \(C_{sj}\) be its earlier guarded
sender bound inclusive of that attempt. Bind at endpoint \(e\) increases
the bound by exactly
\[
 D(s,e,j)=\sum_{\substack{i\text{ before the datagram of }j\\
               s(i)=s,\;e(i)=e}} b_i .
\]
Each earlier datagram to \(e\) gains one possible attempt of size \(b_i\);
all other earlier attempts retain their bound. Within the datagram of
\(j\), an appended binding follows the named receiver, so it contributes
nothing through \(j\) itself. This distinction fails if the membership
change can reorder already bound peers. Actual use can be less if some
attempts are suppressed or old debit expires. Equalities refer to
the conservative prefix-bound update and are tight when every counted
attempt charges before \(j\), with no intervening expiry.

A binding can be committed by this certificate exactly when every
protected sender still fits its bound (and the independently established
non-TX premises still hold). For one capacity and live baseline per
sender, the latest protected attempt has the greatest prefix bound,
so it subsumes the other TX inequalities for that sender. It does not
subsume membership, authority, RX allowance, path or deadline witnesses.
This maximum reduction does not apply independently across sender
budgets, nor across different expiry-aware allowance cuts.

\paragraph{Native distinguishing cases.}
At sender capacities 293 and 233, binding one additional subscriber
to an earlier probe is admissible for the first sender but not the
second. A check on the first alone can break the second report. At
293 for both, both reports and the new subscription fit. When the
changed endpoint affects only later entries, or adds an attempt to a
protected datagram after its old required receiver, the test retains
both reports without deferring the bind. Separately native unregister
removes the needed endpoint binding: there is then plenty of sender
credit but no named receipt. That is not a TX-budget failure and this
guard for new binds is not a repair for membership loss.

\paragraph{Replay and cost.}
Controlled runs begin with two reports that fit under the old binding.
Over 192 rows per arm, the ordinary bind arm never refuses and loses at
least one report in 28 rows. The first-sender-only guard refuses 16
binds and loses a report in another 12. The joint guard defers 28
changes and preserves both reports in all rows. Independent traffic
still delivers at the resumed dispatch. Selected stock Bind runs
unmodified for its baseline; in the guarded fixture a predicate is
inserted just before membership publication, under Bind's own mutex.
No alternative scheduling advantage over the same-information direct
calculation, which agrees in all 576 rows. No claims about measured
contention, service during the pose stall, live registration of a not-yet
pose-ready client, or independent mission acceptance.

The replay scans the pending queue under the existing broker-wide lock;
it does not implement a sparse index. The formula specifies what an
index would retain and the capacity it would invalidate. Maintaining
one under ordinary enqueue, native dequeue, expiry, membership removal
and authority migration, without losing progress to long lock holds,
is a separate implementation obligation. Receiver service costs and
fallback for a rejected request must also be provided in a wider
envelope. The binding guard is a source-qualified extension of the
previous rule, not a claim of a new general transaction principle.
\end{document}
